% Experiment report. Fixed sections: see experiments/TEMPLATE/report.tex and experiments/README.md.
% Build: npm run experiment:report -- TEMPLATE
\documentclass{../polyexp}
\begin{document}
\input{generated}
\ExpHeader

\section{Question and hypothesis}
Stepping through replayed training games of experiment 24 (27 Nov 2026), Jakob marked seven failure modes:
``to the centre should score highest'' (all four first moves of a warrior scored the same, +3.61);
``one of them randomly moved to the edge'';
``a village was discovered by the last unit but stepping towards it now doesn't score higher'';
``this unit stands idle for some reason''; ``does nothing''; ``moves back to the corner''.
The replays show two causes. (1) Every feature describes what the unit can already see; under the fog all directions
look the same, ties are broken at random (edge moves) and from turn 1 every move has gain exactly 0, so the units pass.
(2) Travel distances count whole turns, so a single step toward a target often leaves the distance unchanged.

Jakob: ``add a feature to orient itself on the map: it needs to know whether it's limited by fog of war or far away
from where an enemy city could be.''

\todo{jakob}{Confirm or change the prediction below.}
Behaviour count as in experiment 27: unit-turns that end idle although the unit still had options (experiment 24
linear: 10\% in turns 1--10, 34\% in turns 11--30).
Prediction (Claude's proposal): with the orientation features idle unit-turns in turns 11--30 fall from 34\% to under
20\%, \texttt{tCityOrFog} and \texttt{tFrontier} are among the five most important
features, and the learner minus the champions' mean improves by at least 300 points over the base. If the new
features get weight but play does not change, the missing piece is the reward (experiment 27), not the state.

\section{Mechanism}
Feature set 2 (\texttt{--features 2}, default 1 = experiment 24), in \texttt{features} and \texttt{travel\_fields},
\texttt{engine/unitrl/src/main.rs}:
\begin{itemize}
\item travel to the four targets of experiment 24 (enemy city, village, enemy unit, threatened own city) counted in
  single movement steps on the unit type's move graph (only edges between neighbouring tiles), capped at 16 steps,
  instead of whole turns: every step toward a target changes it;
\item \texttt{tFrontier}, \texttt{noFrontier}: steps to the nearest explored tile next to an unexplored one, and a
  flag when there is none;
\item \texttt{fogNear}: share of unexplored tiles within 2 tiles of the unit;
\item \texttt{exploredShare}: share of the map the learner's player has explored;
\item \texttt{tCityOrFog}: steps to the nearest place an enemy city is or could be: next to a known enemy city, or on
  the fog frontier.
\end{itemize}
All targets are explored tiles, so the features use no hidden information (as in experiment 24, paths may cross
fog). Under feature set 1 the five new features are 0, so experiment 24's runs replay unchanged (checked: iteration
300 replays to the logged mean score, 5776.17).

\section{Setup}
One run, \texttt{rl-unit-orient}: feature set 2, no exploration reward, everything else as experiment 24's linear
run, 300 iterations; base and measurement as in experiment 27. Smoke test (3 iterations): all five new features got
non-zero weights (\texttt{fogNear} $+11.4$, \texttt{tCityOrFog} $-3.9$, \texttt{tFrontier} $-3.1$).
The combination with the exploration reward is a later experiment.

\section{Results}
% WHEN: after the runs. The tables and plots below are generated from the measured data; add two to four sentences that read them.
\ExpResults

On the 200 evaluation games (seeds 6e9.., the same for both), learner minus the champions' mean:
base (experiment 24 linear, iteration 300) $-2322$ (95\% interval $-2834$ to $-1811$), feature set 2 at iteration
300 $-2365$ ($-2874$ to $-1857$). The difference, $-43$, is far inside the noise. Per game: kills 11.9 $\to$ 12.3,
units lost 14.9 $\to$ 17.8, cities captured 0.63 $\to$ 0.72, villages 2.19 $\to$ 2.56. Files:
\texttt{base-measure-200.json}, \texttt{measure-200.json} in this folder.

\section{Behaviour}
% WHEN: after the runs. WHO: both. One count the mechanism should change, plus what was seen in watched games.
Unit-turns that end idle although the unit still had options (12 replayed training games at iteration 300):
17\% in turns 1--10 (experiment 24 linear: 10\%) and 32\% in turns 11--30 (34\%). The units do not idle less.
By permutation importance the value function barely uses the new features: \texttt{noFrontier} ranks 9th of 41,
\texttt{fogNear} 15th, \texttt{tFrontier} 17th, \texttt{tCityOrFog} 30th, \texttt{exploredShare} 35th; the top
three stay \texttt{tVillage}, \texttt{noVillage} and \texttt{turnsLeft}.
\todo{jakob}{What you saw in the live games, if you watched any.}

\section{Interpretation}
% WHEN: after the runs. WHO: both. Why it worked or did not; was the prediction right?
The prediction failed on all three counts (idle unit-turns under 20\%, \texttt{tCityOrFog} and \texttt{tFrontier}
in the top five, +300 points). The features describe the map correctly, but the value function can only use a
feature that predicts the unit's return, and the return contains almost nothing an enemy city or the fog frontier
would predict: cities are captured 0.6 times per game, exploring earns nothing (experiment 27 is the reward side), and
unit deaths cost nothing. Villages do predict returns, so \texttt{tVillage} stays on top and the finer step distance
made the units take a few more villages (+0.37 per game) while losing more units (+2.8). In experiment 24's runs
fighting was also flat from the first 20 iterations to the last: the least-squares fit converges within a few
iterations and greedy play keeps producing the same games, so a new feature that is not tied to a reward cannot
change play.
\todo{jakob}{Your reading: agree, disagree, what is missing?}

\section{Threats to validity}
% WHO: Claude. What could make this result wrong or narrow: panel, map sizes, sample size, interactions with other knobs.
(1) 300 iterations for 41 features against 1000 for experiment 24's final runs; experiment 24's curves were flat
by iteration 300, and the base is measured at iteration 300 too. (2) One run, one seed of the trainer; the idle count
comes from 12 replayed games only. (3) The finer step distance and the five orientation features changed together,
so their separate effects are not known (the step distance may explain the extra villages).

\section{Verdict and next step}
% WHO: Claude proposes by the keep rules (.claude/skills/model-experiment/SKILL.md); Jakob decides anything outside them.
Not kept: no measurable change in score. The code stays (off by default, \texttt{--features 2}). Next: tie the state
to rewards that predict it. A per-unit reward built from the score's own parts (exploration 5 points per tile,
a unit's score value lost when it dies, kills and captures as now), then this feature set again on top of it; and a
way past the early plateau (start V from the champion's unit play, or over-sample captures and kills).

\section{Notes for the paper}
% WHO: Jakob, with Claude's draft. What this experiment contributes: a claim, a figure, a negative result worth a line.
Draft: ``Features that describe where enemies could be did not change play: a value function only uses state
that predicts its return, and with captures this rare the map orientation predicts almost nothing.'' Carried by the
importance ranking (new features 9th--35th of 41) next to the unchanged score.
\todo{jakob}{Is this worth a place in the paper? Where?}

\end{document}
